\documentclass[12pt,letterpaper]{article}

\usepackage[spanish]{babel}
\usepackage[utf8]{inputenc}
\usepackage[T1]{fontenc}
\usepackage{graphicx}
\usepackage{fancyhdr}
\usepackage{geometry}
\usepackage{amsmath}

\geometry{margin=2.5cm}
\setlength{\headheight}{52pt}


\pagestyle{fancy}
\fancyhf{}

\fancyhead[L]{
    \includegraphics[height=1.5cm]{ur_image.png}
}

\fancyhead[R]{
\begin{tabular}{r}
\textbf{Monitoria Álgebra Lineal} \\
Gabriel Fernando Salcedo Zamora \\
\end{tabular}
}

\renewcommand{\headrulewidth}{0.4pt}

\begin{document}

\section*{Ejercicios Rectas y Planos}
\begin{enumerate}
    \item Encuentre una ecuación vectorial, las ecuaciones paramétricas y las simétricas de la recta indicada.
    \begin{enumerate}
        \item Contiene a $(6,10,3)$ y es paralela al vector $(-10,7,9)$
        \item Contiene a $(a,b,c)$ y es paralela al vector $(0,d,e)$
    \end{enumerate}
    \item Determine si las rectas son ortogonales o paralelas.
    \begin{enumerate}
        \item \[
L_1:\displaystyle \frac{x - 3}{2} = \frac{y + 1}{4} = \frac{z - 2}{-1}
\quad
L_2:\displaystyle \frac{x - 3}{5} = \frac{y + 1}{-2} = \frac{z - 3}{2}
\]
        \item \[
L_1:\displaystyle \frac{x - 1}{1} = \frac{y + 3}{2} = \frac{z + 3}{3}
\quad
L_2:\displaystyle \frac{x - 3}{3} = \frac{y - 1}{6} = \frac{z - 8}{9}
\]
    \end{enumerate}
    \item Encuentre la ecuación del plano para cada caso
    \begin{enumerate}
        \item $P=(1,2,3); N=(1,1,0)$
        \item $P=(2,-1,6); N=(3,-1,2)$
    \end{enumerate}
    \item Determine el punto de intersección de la recta: \[
\begin{aligned}
x &= 5 + 2t \\
y &= 1 - 3t \\
z &= 2 + t
\end{aligned}
\] y el plano $2x+3y-4z+5=0$
    \item Determine la ecuación del plano que pase por el punto $(2,4,-3)$ y es paralelo al plano $-2x+4y-5z+6=0$
\end{enumerate}

\newpage

\section*{Respuestas}
\begin{enumerate}
    \item 
    \begin{enumerate}
        \item $(x,y,z) = (6,10,3) + t(-10,7,9)$\[
\begin{aligned}
x &= 6-10t \\
y &= 10+7t \\
z &= 3+9t
\end{aligned}
\]
\[\displaystyle \frac{x - 6}{-10} = \frac{y - 10}{7} = \frac{z - 3}{9}
\]
        \item $(x,y,z) = (a,b,c) + t(0,d,e)$\[
\begin{aligned}
x &= a \\
y &= b+dt \\
z &= c+et
\end{aligned}
\]
\[\displaystyle x=a; \frac{y - b}{d} = \frac{z - c}{e}
\]
    \end{enumerate}
    \item
    \begin{enumerate}
        \item $L_1$ es ortogonal a $L_2$
        \item $L_1$ es paralela a $L_2$
    \end{enumerate}
    \item 
    \begin{enumerate}
        \item $1(x-1)+1(y-2)+0(z-3)=0;\quad x+y-3=0$
        \item $3(x-2)-(y+1)+2(z-6)=0;\quad3x-y+2z-19=0$
    \end{enumerate}
    \item El punto es: $(\frac{65}{9}, -\frac{7}{3}, \frac{28}{9})$
    \item La ecuación del plano es: $-2x+4y-5z-27=0$
\end{enumerate}
\end{document}
